% Part: normal-modal-logic
% Chapter: tableaux
% Section: countermodels

\documentclass[../../../include/open-logic-section]{subfiles}

\begin{document}

\olfileid[zh]{nml}{tab}{cou}

\olsection{来自\usetoken{P}{tableau}的反模型}

\begin{explain}
  完全性定理的证明不仅表明，若 $\Entails !A$，则 $\Proves !A$；它还给出了一个方法：若 $\Entails/ A$，则为~$!A$ 构造反模型。就~$\Log{K}$ 而言，这一方法构成一个\emph{判定程序}。设 $\Entails/ !A$。则 \olref[cpl]{prop:complete-tableau} 的证明给出一个构造完全!!{tableau}的方法。事实上，这一方法总会终止。$\Log{K}$~的命题规则只添加复杂度更低的带前缀!!{formula}；即，对任何带符号公式 $\sFmla{S}{!A}[\sigma]$，每条命题规则在一支上只需应用一次。新的前缀只由
  \iftag{prvBox}{$\TRule{\False}{\Box}$}{}\iftag{notprvBox,notprvDiamond}{}{
    和 }\iftag{prvDiamond}{$\TRule{\True}{\Diamond}$}{}
  \iftag{notprvBox,notprvDiamond}{规则}{规则}生成，而且这些规则也只需应用一次（并产生一个新的前缀）。
  \iftag{prvBox}{$\TRule{\True}{\Box}$}{}\iftag{notprvBox,notprvDiamond}{}{
    和 }\iftag{prvDiamond}{$\TRule{\False}{\Diamond}$}{}
  \iftag{notprvBox,notprvDiamond}{}{}可能需要多次应用，但对每个前缀只需应用一次，而且只会生成有穷多个新的前缀。因此，这一构造在有限多个阶段之后，要么产生一个闭的支，要么产生一个完全的支。

  一旦构造出具有一个开放且完全的支的!!{tableau}，\olref[cpl]{thm:tableau-completeness} 的证明就给出一个满足原始带前缀!!{formula}集合的显式模型。因此，我们得到的不只是：若 $\Gamma \Entails !A$，则存在闭的!!{tableau}，且 $\Gamma \Proves !A$；若以适当的方式寻找闭的!!{tableau}，而最终得到一个「完全」的!!{tableau}，那么我们不仅会知道 $\Gamma \Entails/ !A$，而且实际上能够构造一个反模型。
\end{explain}

\iftag{prvBox}{
\begin{ex}
  我们知道 $\Proves/ \Box(p \lor q) \lif (\Box p \lor \Box
  q)$。表列的构造从如下开始：
  \begin{oltableau}
    [\pFmla{\False}{\Box(p \lor q) \lif (\Box p \lor \Box q)}{1},
      just = \TAss, checked
      [\pFmla{\True}{\Box(p \lor q)}{1},
        just = {\TRule{\False}{\lif}[1]},
        [\pFmla{\False}{\Box p \lor \Box q}{1},
          just = {\TRule{\False}{\lif}[1]}, checked
          [\pFmla{\False}{\Box p}{1},
            just = {\TRule{\False}{\lor}[3]}, checked
            [\pFmla{\False}{\Box q}{1},
              just = {\TRule{\False}{\lor}[3]}, checked
              [\pFmla{\False}{p}{1.1},
                just = {\TRule{\False}{\Box}[4]}, checked
                [\pFmla{\False}{q}{1.2},
                  just = {\TRule{\False}{\Box}[5]}, checked
                ]
              ]
            ]
          ]
        ]
      ]
    ]
  \end{oltableau}
  当然，!!{tableau}还没有完成。下一步，我们考虑唯一没有对勾的行：第~$2$~行的带前缀!!{formula}
  $\sFmla{\True}{\Box(p \lor q)}[1]$。按照闭!!{tableau}的构造方法，需要对支上使用的每个前缀应用 $\TRule{\True}{\Box}$ 规则，即对 $1.1$ 和 $1.2$ 两者都应用：
  \begin{oltableau}
    [\pFmla{\False}{\Box(p \lor q) \lif (\Box p \lor \Box q)}{1},
      just = \TAss, checked
      [\pFmla{\True}{\Box(p \lor q)}{1},
        just = {\TRule{\False}{\lif}[1]},
        [\pFmla{\False}{\Box p \lor \Box q}{1},
          just = {\TRule{\False}{\lif}[1]}, checked
          [\pFmla{\False}{\Box p}{1},
            just = {\TRule{\False}{\lor}[3]}, checked
            [\pFmla{\False}{\Box q}{1},
              just = {\TRule{\False}{\lor}[3]}, checked
              [\pFmla{\False}{p}{1.1},
                just = {\TRule{\False}{\Box}[4]}, checked
                [\pFmla{\False}{q}{1.2},
                  just = {\TRule{\False}{\Box}[5]}, checked
                  [\pFmla{\True}{p \lor q}{1.1},
                    just = {\TRule{\True}{\Box}[2]}
                    [\pFmla{\True}{p \lor q}{1.2},
                      just = {\TRule{\True}{\Box}[2]}
                    ]
                  ]
                ]
              ]
            ]
          ]
        ]
      ]
    ]
  \end{oltableau}
  现在第 2、8、9 行没有对勾。但没有添加新的前缀，因此我们在由此产生的所有支上（只要它们尚未闭合）对第~8~行和第~9~行应用 $\TRule{\True}{\lor}$：
  \begin{oltableau}
    [\pFmla{\False}{\Box(p \lor q) \lif (\Box p \lor \Box q)}{1},
      just = \TAss, checked
      [\pFmla{\True}{\Box(p \lor q)}{1},
        just = {\TRule{\False}{\lif}[1]}, checked
        [\pFmla{\False}{\Box p \lor \Box q}{1},
          just = {\TRule{\False}{\lif}[1]}, checked
          [\pFmla{\False}{\Box p}{1},
            just = {\TRule{\False}{\lor}[3]}, checked
            [\pFmla{\False}{\Box q}{1},
              just = {\TRule{\False}{\lor}[3]}, checked
              [\pFmla{\False}{p}{1.1},
                just = {\TRule{\False}{\Box}[4]}, checked
                [\pFmla{\False}{q}{1.2},
                  just = {\TRule{\False}{\Box}[5]}, checked
                  [\pFmla{\True}{p \lor q}{1.1},
                    just = {\TRule{\True}{\Box}[2]}, checked
                    [\pFmla{\True}{p \lor q}{1.2},
                      just = {\TRule{\True}{\Box}[2]}, checked
                      [\pFmla{\True}{p}{1.1},
                        just = {\TRule{\True}{\lor}[8]}, checked, close
                      ]
                      [\pFmla{\True}{q}{1.1},
                        just = {\TRule{\True}{\lor}[8]}, checked
                        [\pFmla{\True}{p}{1.2},
                          just = {\TRule{\True}{\lor}[9]}, checked]
                        [\pFmla{\True}{q}{1.2},
                          just = {\TRule{\True}{\lor}[9]}, checked, close]
                      ]
                    ]
                  ]
                ]
              ]
            ]
          ]
        ]
      ]
    ]
  \end{oltableau}
  还剩下一个开放的支，并且它是完全的。我们从它定义模型：世界 $W = \{1, 1.1, 1.2\}$（开放支上出现的唯一一些前缀）、可及关系
  $R = \{\tuple{1, 1.1}, \tuple{1, 1.2}\}$，以及指派 $V(p) =
  \{1.2\}$（因为第~11~行包含 $\sFmla{\True}{p}[1.2]$）和
  $V(q) = \{1.1\}$（因为第~10~行包含
  $\sFmla{\True}{q}[1.1]$）。该模型如图
  \olref{fig:counter-Box} 所示，你可以验证它是
  $\Box(p \lor q) \lif (\Box p \lor \Box q)$ 的一个反模型。
  \begin{figure}
  \begin{center}
    \begin{tikzpicture}[modal]
      \node[world] (w1) [label={[align=right]right:\mFalse{p}\\ \mFalse{q}}]{$1$};
      \node[world] (w2) [label={[align=right]right:\mFalse{p}\\ \mTrue{q}},
        above left=of w1]{$1.1$};
      \node[world] (w3) [label={[align=right]right:\mTrue{p}\\ \mFalse{q}},
        above right=of w1] {$1.2$};
      \draw[->] (w1) to (w2);
      \draw[->] (w1) to (w3);
    \end{tikzpicture}
  \end{center}
  \caption{$\Box(p \lor q) \lif (\Box p \lor \Box
    q)$ 的一个反模型。}
  \ollabel{fig:counter-Box}
\end{figure}
\end{ex}
}
{
\begin{ex}
  我们知道 $\Proves/ (\Diamond p \land \Diamond q) \lif \Diamond(p
  \land q)$。表列的构造从如下开始：
  \begin{oltableau}
    [\pFmla{\False}{(\Diamond p \land \Diamond q) \lif \Diamond(p \land q)}{1},
      just = \TAss, checked
      [\pFmla{\True}{\Diamond p \land \Diamond q}{1},
        just = {\TRule{\False}{\lif}[1]}, checked
        [\pFmla{\False}{\Diamond(p \land q)}{1},
          just = {\TRule{\False}{\lif}[1]}
          [\pFmla{\True}{\Diamond p}{1},
            just = {\TRule{\True}{\land}[2]}, checked
            [\pFmla{\True}{\Diamond q}{1},
              just = {\TRule{\True}{\land}[2]}, checked
              [\pFmla{\True}{p}{1.1},
                just = {\TRule{\True}{\Diamond}[4]}, checked
                [\pFmla{\True}{q}{1.2},
                  just = {\TRule{\True}{\Diamond}[5]}, checked
                ]
              ]
            ]
          ]
        ]
      ]
    ]
  \end{oltableau}
  当然，!!{tableau}还没有完成。下一步，我们考虑唯一没有对勾的行：第~$3$~行的带前缀!!{formula}
  $\sFmla{\True}{\Diamond(p \land q)}[1]$。闭
  表列的构造要求对支上使用的每个前缀应用 $\TRule{\True}{\Diamond}$ 规则，即对 $1.1$ 和 $1.2$ 两者都应用：
  \begin{oltableau}
    [\pFmla{\False}{\Diamond(p \land q) \lif (\Diamond p \land \Diamond q)}{1},
      just = \TAss, checked
      [\pFmla{\True}{\Diamond p \land \Diamond q}{1},
        just = {\TRule{\False}{\lif}[1]}, checked
        [\pFmla{\False}{\Diamond (p \land q)}{1},
          just = {\TRule{\False}{\lif}[1]}
          [\pFmla{\True}{\Diamond p}{1},
            just = {\TRule{\True}{\land}[2]}, checked
            [\pFmla{\True}{\Diamond q}{1},
              just = {\TRule{\True}{\land}[2]}, checked
              [\pFmla{\True}{p}{1.1},
                just = {\TRule{\True}{\Diamond}[4]}, checked
                [\pFmla{\True}{q}{1.2},
                  just = {\TRule{\True}{\Diamond}[5]}, checked
                  [\pFmla{\False}{p \land q}{1.1},
                    just = {\TRule{\False}{\Diamond}[3]}
                    [\pFmla{\False}{p \land q}{1.2},
                      just = {\TRule{\False}{\Diamond}[3]}
                    ]
                  ]
                ]
              ]
            ]
          ]
        ]
      ]
    ]
  \end{oltableau}
    现在第 3、8、9 行没有对勾。但没有添加新的前缀，因此我们在由此产生的所有支上（只要它们尚未闭合）对第~8~行和第~9~行应用 $\TRule{\False}{\land}$：  
    \begin{oltableau}
    [\pFmla{\False}{(\Diamond p \land \Diamond q) \lif \Diamond(p \land q)}{1},
      just = \TAss, checked
      [\pFmla{\True}{\Diamond p \land \Diamond q}{1},
        just = {\TRule{\False}{\lif}[1]}, checked
        [\pFmla{\False}{\Diamond(p \land q)}{1},
          just = {\TRule{\False}{\lif}[1]}
          [\pFmla{\True}{\Diamond p}{1},
            just = {\TRule{\True}{\land}[2]}, checked
            [\pFmla{\True}{\Diamond q}{1},
              just = {\TRule{\True}{\land}[2]}, checked
              [\pFmla{\True}{p}{1.1},
                just = {\TRule{\True}{\Diamond}[4]}, checked
                [\pFmla{\True}{q}{1.2},
                  just = {\TRule{\True}{\Diamond}[5]}, checked
                  [\pFmla{\False}{p \land q}{1.1},
                    just = {\TRule{\False}{\Diamond}[3]}, checked
                    [\pFmla{\False}{p \land q}{1.2},
                      just = {\TRule{\False}{\Diamond}[3]}, checked
                      [\pFmla{\False}{p}{1.1},
                        just = {\TRule{\False}{\land}[8]}, checked, close
                      ]
                      [\pFmla{\False}{q}{1.1},
                        just = {\TRule{\False}{\land}[8]}, checked
                        [\pFmla{\False}{p}{1.2},
                          just = {\TRule{\False}{\land}[9]}, checked]
                        [\pFmla{\False}{q}{1.2},
                          just = {\TRule{\False}{\land}[9]}, checked, close]
                      ]
                    ]
                  ]
                ]
              ]
            ]
          ]
        ]
      ]
    ]
  \end{oltableau}
  还剩下一个开放的支，并且它是完全的。我们由此定义模型：世界 $W = \{1, 1.1, 1.2\}$（开放支上出现的全部前缀）、可及关系
  $R = \{\tuple{1, 1.1}, \tuple{1, 1.2}\}$，以及指派 $V(p) =
  \{1.1\}$（因为第~6~行包含 $\sFmla{\True}{p}[1.1]$）和
  $V(q) = \{1.2\}$（因为第~7~行包含
  $\sFmla{\True}{q}[1.1]$）。该模型如图
  \olref{fig:counter-Diamond} 所示，你可以验证它是
  $(\Diamond p \land \Diamond q) \lif \Diamond (p \land q)$ 的一个反模型。
  \begin{figure}
  \begin{center}
    \begin{tikzpicture}[modal]
      \node[world] (w1) [label={[align=right]right:\mFalse{p}\\ \mFalse{q}}]{$1$};
      \node[world] (w2) [label={[align=right]right:\mTrue{p}\\ \mFalse{q}},
        above left=of w1]{$1.1$};
      \node[world] (w3) [label={[align=right]right:\mFalse{p}\\ \mTrue{q}},
        above right=of w1] {$1.2$};
      \draw[->] (w1) to (w2);
      \draw[->] (w1) to (w3);
    \end{tikzpicture}
  \end{center}
  \caption{$(\Diamond p \land \Diamond q) \lif
    \Diamond (p \land q)$ 的一个反模型。}
  \ollabel{fig:counter-Diamond}
\end{figure}
\end{ex}
}

\end{document}
